\section{Discussion}
\label{sec:discussion}

The central claim of this paper is that asking a human should be treated as part of the robot's belief dynamics, not as an external recovery behavior. This matters because the usefulness of a query depends on both the robot's current posterior and the task consequences of committing the wrong symbolic state. A query about a visually ambiguous tomato is valuable if the answer changes whether the robot harvests, discards, or ignores it. The same query may be unnecessary if the tomato is irrelevant to the current goal or if another action can resolve the uncertainty autonomously.

The frontier belief representation makes this connection explicit. By constructing uncertainty around states reachable from the current knowledge state and selected action, the robot avoids asking about facts that do not affect near-term planning. The entropy-based confidence score determines when the robot lacks enough information to commit a state, and the expected-entropy query rule determines which symbolic fact would most reduce the ambiguity. In this formulation, the knowledge base is not simply a database of perceptual outputs. It is a guarded state estimator: local perceptual facts become planning knowledge only after the belief is sufficiently concentrated.

The system evaluation supports this interpretation. The sharp performance change around $\tau=0.7$ suggests that the confidence threshold functions as a meaningful guard between ambiguous belief and committed knowledge. The policy comparison further shows that query timing cannot be replaced by query quantity alone: random feedback can spend interactions without stabilizing the belief states that matter for planning, while always querying incurs unnecessary interaction cost.

The user study should be interpreted as secondary evidence rather than the core contribution. Its results suggest that state-level queries can make the robot's uncertainty more legible and can reduce workload relative to action-level query baselines. However, the paper's main claim is not that one interaction condition universally improves all human factors measures. Rather, the human results are consistent with the system design principle that people should refine uncertain state information while the robot remains responsible for belief update and action selection.

Several limitations remain. First, the framework assumes that human answers are reliable or at least more reliable than the robot's uncertain observation. Future work should model noisy human feedback as another observation source rather than as a hard filter. Second, the current query form is binary and predicate-level. More expressive questions may reduce the number of interactions but increase user burden. Third, the planner uses a local frontier belief rather than a full global posterior. This improves tractability and preserves symbolic action relevance, but it may miss dependencies outside the immediate action-conditioned frontier. Finally, the current implementation relies on domain-specific frontier construction; learning or automatically generating these frontier models would make the framework easier to transfer across tasks.
